% fig_sounding.tex -- one 802.11 sounding as a two-lifeline sequence diagram.
%
% Single-column figure (\columnwidth = 241.15pt = 84.75mm wide, <= 1.5in tall).
% Requires in the document preamble:
%     \usepackage{tikz}
%     \usetikzlibrary{arrows.meta,positioning,calc}
%
% Palette: Okabe-Ito (access point blue, client green, feedback frame orange,
% outsider vermilion), as in the rest of the paper.
%
% Layout (time runs downward):
%   1. NDP announcement             access point -> client
%   2. NDP (known training symbols) access point -> client, note at the client
%   3. Compressed Beamforming Action frame, client -> access point, thick orange,
%      with a branch from its midpoint to an outsider's mast at the far right.
%
\begin{tikzpicture}[
    x=1mm, y=1mm,
    font=\footnotesize,
    >={Stealth[length=1.4mm,width=1.1mm]},
    inner sep=0pt, outer sep=0pt,
]

% ===== Colour palette (Okabe-Ito) =====
\definecolor{cBlue}{HTML}{0072B2}
\definecolor{cGreen}{HTML}{009E73}
\definecolor{cOrange}{HTML}{E69F00}
\definecolor{cRed}{HTML}{D55E00}
\definecolor{cGray}{HTML}{555555}
\definecolor{fillGw}{HTML}{DCEEF7}
\definecolor{fillCl}{HTML}{D9F0E8}

% Lifeline abscissae; labels between the lifelines are centred on \xm.
\def\xg{9}
\def\xc{50}
\def\xm{29.5}
\def\xo{77}
\def\xr{84.5}

% ----- Header boxes (lifeline tops) -----
\node[draw=cBlue, fill=fillGw, line width=0.4pt, rounded corners=0.7mm,
      inner sep=0.7mm, minimum width=12mm, minimum height=4mm, anchor=south]
      (gw) at (\xg,31) {access point};
\node[draw=cGreen, fill=fillCl, line width=0.4pt, rounded corners=0.7mm,
      inner sep=0.7mm, minimum width=12mm, minimum height=4mm, anchor=south]
      (cl) at (\xc,31) {client};

% ----- Lifelines -----
\draw[cGray, line width=0.3pt, dash pattern=on 0.8mm off 0.5mm] (gw.south) -- (\xg,4);
\draw[cGray, line width=0.3pt, dash pattern=on 0.8mm off 0.5mm] (cl.south) -- (\xc,4);

% ----- Time arrow on far left -----
\draw[->, cGray, line width=0.3pt] (1.5,29) -- (1.5,6);
\node[cGray, rotate=90, font=\scriptsize] at (3.1,17.5) {time};

% ===== 1. NDP announcement: access point -> client =====
\draw[->, cBlue, line width=0.5pt] (\xg+0.6,27.5) -- (\xc-0.6,27.5);
\node[anchor=south] at (\xm,27.9) {NDP announcement};

% ===== 2. NDP with known training symbols: access point -> client =====
\draw[->, cBlue, line width=0.5pt] (\xg+0.6,21.5) -- (\xc-0.6,21.5);
\node[anchor=south] at (\xm,21.9) {NDP (known training symbols)};
\node[anchor=west, cGray, font=\scriptsize\itshape, align=left] at (\xc+1.8,21.5)
     {client estimates\\the channel};

% ===== 3. Compressed Beamforming Action frame: client -> access point =====
\draw[->, cOrange, line width=1.0pt] (\xc-0.6,8.5) -- (\xg+0.6,8.5);
\node[anchor=south, align=center, fill=white, inner sep=0.3mm] at (\xm,9.1)
     {Compressed Beamforming Report:\\SNR per stream, Givens angles\\per group, 2 to 9 bits, plaintext};

% ----- Branch from the midpoint of the feedback frame to the outsider -----
\draw[->, cRed, line width=0.5pt, dash pattern=on 1mm off 0.6mm, rounded corners=1.5mm]
    (\xm,8.5) |- (\xo-2.2,5.5);

% ----- Outsider mast (far right) -----
\draw[cRed, line width=0.6pt] (\xo-1.5,5.5) -- (\xo+1.5,5.5);    % footing
\draw[cRed, line width=0.6pt] (\xo,5.5) -- (\xo,11);               % mast
\fill[cRed] (\xo,11) circle (0.5);                                 % tip
\draw[cRed, line width=0.4pt] (\xo-1.6,12.2) arc (120:60:3.2);     % wave 1
\draw[cRed, line width=0.4pt] (\xo-2.6,13.4) arc (120:60:5.2);     % wave 2
\node[cRed, anchor=south east] at (\xr,14.2) {outsider in range};
\node[cGray, font=\scriptsize\itshape, anchor=north east, align=right] at (\xr,4.7)
     {not covered by\\Management Frame Protection};

\end{tikzpicture}
